\documentclass[11pt, a4paper]{article}

% Gemeinsame Style-Definitionen (TEXINPUTS muss templates/ enthalten — siehe scripts/build_tex.py)
% bfw-style.tex — gemeinsame Style-Definitionen für alle Handouts/Aufgabenhefte
% Voraussetzung: Kompilation mit xelatex (wegen fontspec)

\usepackage[ngerman]{babel}
\usepackage{geometry}
\geometry{a4paper, top=2.2cm, bottom=2.2cm, left=2.3cm, right=2.3cm}

\usepackage{fontspec}
% Fallback-Kette: Source Sans Pro -> Calibri -> Segoe UI -> Latin Modern Sans
% (Latin Modern ist immer mit MiKTeX vorhanden)
\IfFontExistsTF{Source Sans Pro}{%
  \setmainfont{Source Sans Pro}[Ligatures=TeX]%
}{\IfFontExistsTF{Calibri}{%
  \setmainfont{Calibri}[Ligatures=TeX]%
}{\IfFontExistsTF{Segoe UI}{%
  \setmainfont{Segoe UI}[Ligatures=TeX]%
}{%
  \setmainfont{Latin Modern Sans}[Ligatures=TeX]%
}}}
\IfFontExistsTF{Source Code Pro}{%
  \setmonofont{Source Code Pro}[Scale=0.92]%
}{\IfFontExistsTF{Consolas}{%
  \setmonofont{Consolas}[Scale=0.92]%
}{%
  \setmonofont{Latin Modern Mono}[Scale=0.92]%
}}

\usepackage{xcolor}
% Farbpalette: hell, freundlich, modern
\definecolor{bfwPrimary}{HTML}{0EA5A4}    % Petrol/Türkis - Primär
\definecolor{bfwPrimaryDark}{HTML}{0F766E} % dunkler Petrol für Headings
\definecolor{bfwAccent}{HTML}{F59E0B}     % Amber - Akzent
\definecolor{bfwSuccess}{HTML}{10B981}    % Grün - Erfolg / Lösung
\definecolor{bfwWarn}{HTML}{F97316}       % Orange - Achtung
\definecolor{bfwInk}{HTML}{0F172A}        % fast Schwarz
\definecolor{bfwInkSoft}{HTML}{475569}    % gedämpftes Grau
\definecolor{bfwBgSoft}{HTML}{F8FAFC}     % off-white
\definecolor{bfwBgTeal}{HTML}{ECFEFF}     % helles Türkis
\definecolor{bfwBgAmber}{HTML}{FEF3C7}    % helles Gelb
\definecolor{bfwBgRose}{HTML}{FFE4E6}     % helles Rosé für Eselsbrücken

\usepackage[colorlinks=true, linkcolor=bfwPrimaryDark, urlcolor=bfwPrimaryDark]{hyperref}
\usepackage{enumitem}
\setlist[itemize]{leftmargin=*, itemsep=2pt, topsep=2pt}
\setlist[enumerate]{leftmargin=*, itemsep=2pt, topsep=2pt}

\usepackage[most]{tcolorbox}
\usepackage{booktabs}
\usepackage{tikz}
\usetikzlibrary{decorations.pathmorphing, positioning, shapes.geometric, arrows.meta}

% Merkkästchen — wichtige Definition / Kernpunkt
\newtcolorbox{merksatz}[1][]{%
  enhanced, breakable,
  colback=bfwBgTeal,
  colframe=bfwPrimary,
  boxrule=0.6pt,
  arc=4pt,
  left=10pt, right=10pt, top=8pt, bottom=8pt,
  fonttitle=\bfseries\color{bfwPrimaryDark},
  title={\faLightbulb\ Merksatz},
  #1
}

% Eselsbrücke — Mnemoniks
\newtcolorbox{eselsbruecke}[1][]{%
  enhanced, breakable,
  colback=bfwBgRose,
  colframe=bfwAccent,
  boxrule=0.6pt,
  arc=4pt,
  left=10pt, right=10pt, top=8pt, bottom=8pt,
  fonttitle=\bfseries\color{bfwWarn},
  title={Eselsbrücke},
  #1
}

% Beispiel-Box
\newtcolorbox{beispiel}[1][]{%
  enhanced, breakable,
  colback=bfwBgSoft,
  colframe=bfwInkSoft,
  boxrule=0.4pt,
  arc=3pt,
  left=10pt, right=10pt, top=8pt, bottom=8pt,
  fonttitle=\bfseries\color{bfwInk},
  title={Beispiel},
  #1
}

% Lösungs-Box (für Aufgabenhefte)
\newtcolorbox{loesung}[1][]{%
  enhanced, breakable,
  colback=bfwBgAmber!40,
  colframe=bfwSuccess,
  boxrule=0.6pt,
  arc=3pt,
  left=10pt, right=10pt, top=8pt, bottom=8pt,
  fonttitle=\bfseries\color{bfwSuccess!70!black},
  title={Lösung},
  #1
}

% Glossar-Eintrag
\newcommand{\glossareintrag}[2]{%
  \par\noindent\textbf{\color{bfwPrimaryDark}#1} \enspace #2\par\smallskip
}

% Section-Stil — etwas farbiger als Standard
\usepackage{titlesec}
\titleformat{\section}
  {\Large\bfseries\color{bfwPrimaryDark}}
  {\thesection}{0.6em}{}
\titleformat{\subsection}
  {\large\bfseries\color{bfwPrimary}}
  {\thesubsection}{0.5em}{}
\titleformat{\subsubsection}
  {\normalsize\bfseries\color{bfwInk}}
  {\thesubsubsection}{0.4em}{}

% TikZ-Stil "skizzenhaft" — etwas weicher als Rough.js, aber stabil
\tikzset{
  roughbox/.style = {
    draw=bfwPrimaryDark, thick, fill=bfwBgTeal,
    rounded corners=4pt, inner sep=6pt
  },
  roughaccent/.style = {
    draw=bfwAccent, thick, fill=bfwBgAmber,
    rounded corners=4pt, inner sep=6pt
  },
  roughline/.style = {
    draw=bfwInkSoft, thick, -{Stealth[length=2.5mm]}
  }
}

% Bullet-Symbole (bewusst kein FontAwesome — vermeidet Font-Installations-Probleme)
\newcommand{\faLightbulb}{$\bullet$}
\newcommand{\faBookmark}{$\bullet$}

% Header/Footer
\usepackage{fancyhdr}
\pagestyle{fancy}
\fancyhf{}
\renewcommand{\headrulewidth}{0pt}
\fancyfoot[L]{\small\color{bfwInkSoft}\thepage}
\fancyfoot[R]{\small\color{bfwInkSoft} BFW M-064 Beschaffung im Büromanagement}


\title{\vspace*{-1.5cm}\Huge\bfseries\color{bfwPrimaryDark}Session~1: Rechtliche Grundlagen der Berufsausbildung\\[0.4em]
  \large\color{bfwInkSoft}\textnormal{Vom Ausbildungsvertrag bis zur Mitbestimmung}}
\author{\textsf{Lernfeld 1 --- Die eigene Rolle im Betrieb mitgestalten \textperiodcentered{} \small\copyright\ Can Siebert\,\textperiodcentered\,2026}}
\date{}

\begin{document}
\maketitle
\thispagestyle{empty}

\vspace{1em}
\begin{merksatz}[title={\bfseries Worum geht es in dieser Session?}]
Als Umschüler:in stehen Sie in einem Berufsausbildungsverhältnis --- einem Rechtsverhältnis
mit klaren Regeln, Rechten und Pflichten. Wer diese Regeln kennt, kann die eigene Rolle im
Betrieb sicher gestalten. In diesem Handout lernen Sie, wie das duale System funktioniert,
welche Gesetze die Ausbildung ordnen (vor allem das Berufsbildungsgesetz), was im
Berufsausbildungsvertrag stehen muss und welche Pflichten beide Seiten haben. Außerdem
geht es um den besonderen Schutz junger Menschen (Jugendarbeitsschutzgesetz) und um die
Mitbestimmung durch Betriebsrat und Jugend- und Auszubildendenvertretung. Das ist die
Grundlage für Lernfeld~1 und für die IHK-Prüfung.
\end{merksatz}

\tableofcontents
\vspace{1em}
\hrule
\vspace{1em}

\section{Das duale System und die eigene Rolle}

Die Berufsausbildung in Deutschland wird als \emph{duale Ausbildung} bezeichnet, weil sie
parallel an zwei \emph{Lernorten} stattfindet: im \textbf{Betrieb} und in der
\textbf{Berufsschule}. Der Betrieb vermittelt die praktische berufliche Handlungsfähigkeit,
die Berufsschule die fachtheoretischen und allgemeinbildenden Kenntnisse. Rechtliche
Grundlage des betrieblichen Teils ist das \emph{Berufsbildungsgesetz (BBiG)}, das 1969 in
Kraft trat und zuletzt mit Wirkung ab dem 01.\ Januar 2020 novelliert wurde.

Seit dem 1.\ August 2014 sind die drei früheren Büroberufe zum Beruf
\emph{Kaufmann/-frau für Büromanagement} zusammengefasst. Die Ausbildung dauert in der
Regel drei Jahre und findet in den Bereichen Industrie und Handel, Handwerk oder
Öffentlicher Dienst statt.

An einem Ausbildungsverhältnis sind drei Rollen beteiligt, die man sauber unterscheiden
muss:

\begin{itemize}
  \item \textbf{Ausbildende:r} --- der Vertragspartner, also das Unternehmen bzw.\ die
    Inhaberin oder der Inhaber. Wer andere zur Berufsausbildung einstellt, hat mit ihnen
    einen Berufsausbildungsvertrag abzuschließen (\S~10 BBiG).
  \item \textbf{Ausbilder:in} --- die Person, die die Ausbildung im Betrieb praktisch
    durchführt. Die/der Ausbildende muss selbst ausbilden oder ausdrücklich eine:n
    Ausbilder:in damit beauftragen.
  \item \textbf{Auszubildende:r} --- die Person, die ausgebildet wird und sich um den Erwerb
    der beruflichen Handlungsfähigkeit bemüht.
\end{itemize}

\begin{merksatz}
Verwechseln Sie nicht \emph{Ausbildende:r} und \emph{Ausbilder:in}: Die/der \textbf{Ausbildende}
ist der Vertragspartner (das Unternehmen), die/der \textbf{Ausbilder:in} ist die Person, die
Ihnen den Beruf beibringt.
\end{merksatz}

\section{Rechtsgrundlagen und der Berufsausbildungsvertrag}

\subsection{Von der Ausbildungsordnung zum betrieblichen Ausbildungsplan}

Für jeden anerkannten Ausbildungsberuf wird nach \S~5 BBiG eine \emph{Ausbildungsordnung}
erlassen. Sie enthält u.\,a.\ die Bezeichnung des Berufs, die Ausbildungsdauer, das
\emph{Ausbildungsberufsbild} (die zu vermittelnden Fertigkeiten, Kenntnisse und Fähigkeiten),
den \emph{Ausbildungsrahmenplan} sowie die Prüfungsanforderungen.

Der \textbf{Ausbildungsrahmenplan} gibt eine grobe sachliche und zeitliche Gliederung vor ---
bundeseinheitlich für alle Betriebe. Er dient der/dem Ausbildenden als Vorlage für den
\textbf{betrieblichen Ausbildungsplan}, den das Unternehmen individuell erstellt und dem
Berufsausbildungsvertrag als Anlage beifügt.

\begin{eselsbruecke}
Denken Sie an einen \textbf{Trichter}: \emph{BBiG} (Gesetz) $\rightarrow$
\emph{Ausbildungsordnung} (je Beruf) $\rightarrow$ \emph{Ausbildungsrahmenplan}
(bundesweiter Fahrplan) $\rightarrow$ \emph{betrieblicher Ausbildungsplan} (Ihr Betrieb).
Vom Allgemeinen zum Konkreten.
\end{eselsbruecke}

\subsection{Abschluss, Form und Eintragung}

Nach \S~10 BBiG hat, wer andere zur Berufsausbildung einstellt, mit ihnen einen
\emph{Berufsausbildungsvertrag} abzuschließen. Der Vertrag wird von der/dem Ausbildenden und
der/dem Auszubildenden unterschrieben; ist die/der Auszubildende noch keine 18~Jahre alt,
müssen zusätzlich die gesetzlichen Vertreter (i.\,d.\,R.\ beide Elternteile) unterzeichnen.
Nach \S~11 Abs.~1 BBiG ist der wesentliche Vertragsinhalt \emph{schriftlich} niederzulegen;
die \emph{elektronische Form ist ausgeschlossen}. Eine Ausfertigung ist der/dem Auszubildenden
und den gesetzlichen Vertretern unverzüglich auszuhändigen.

Das Berufsausbildungsverhältnis wird bei der \emph{zuständigen Stelle} eingetragen --- für
kaufmännische Berufe ist das die \textbf{Industrie- und Handelskammer (IHK)}. Sie führt das
Verzeichnis der Berufsausbildungsverhältnisse; Änderungen des wesentlichen Vertragsinhalts
sind ihr unverzüglich anzuzeigen.

\subsection{Die neun Pflichtinhalte des Vertrages}

Nach \S~11 BBiG müssen folgende neun Punkte in den Vertrag aufgenommen werden:

\begin{enumerate}
  \item Art, sachliche und zeitliche Gliederung sowie Ziel der Berufsausbildung.
  \item Beginn und Dauer der Berufsausbildung.
  \item Ausbildungsmaßnahmen außerhalb der Ausbildungsstätte.
  \item Dauer der regelmäßigen täglichen Ausbildungszeit.
  \item Dauer der Probezeit.
  \item Zahlung und Höhe der Vergütung.
  \item Dauer des Urlaubs.
  \item Voraussetzungen, unter denen der Vertrag gekündigt werden kann.
  \item Hinweis auf anzuwendende Tarifverträge, Betriebs- oder Dienstvereinbarungen.
\end{enumerate}

\subsection{Probezeit, Vergütung und Urlaub}

Das Berufsausbildungsverhältnis beginnt mit der \textbf{Probezeit}. Sie muss mindestens
\emph{einen Monat} und darf höchstens \emph{vier Monate} betragen (\S~20 BBiG).

Nach \S~17 BBiG ist eine \textbf{angemessene Vergütung} zu zahlen, deren Höhe im Vertrag
angegeben sein muss. Die Vergütung muss mindestens \emph{jährlich steigen}. Für Verträge ab
dem 01.01.2020 gilt zusätzlich eine gesetzliche Mindestausbildungsvergütung; bei
tarifgebundenen Arbeitgebern gilt die tarifvertragliche Vergütung.

Auch die Dauer des \textbf{Urlaubs} muss im Vertrag stehen. Für erwachsene Auszubildende gilt
das \emph{Bundesurlaubsgesetz (BUrlG)}: mindestens 24~Werktage jährlich (\S~3 BUrlG). Für
Jugendliche unter 18~Jahren gilt der gestaffelte Mindesturlaub des
Jugendarbeitsschutzgesetzes (siehe Abschnitt~4).

\begin{beispiel}
Ein 17-jähriger Auszubildender schließt einen Vertrag mit einer Probezeit von drei Monaten.
Das ist zulässig, denn die Probezeit liegt im gesetzlichen Rahmen von einem bis vier Monaten.
Eine Probezeit von fünf Monaten wäre unwirksam.
\end{beispiel}

\section{Pflichten, Beendigung und Zeugnis}

\subsection{Pflichten der Auszubildenden (\S~13 BBiG)}

Auszubildende haben sich zu bemühen, die berufliche Handlungsfähigkeit zu erwerben.
Insbesondere sind sie verpflichtet, aufgetragene Aufgaben \emph{sorgfältig} auszuführen, an
Ausbildungsmaßnahmen teilzunehmen, \emph{Weisungen} zu folgen, die \emph{Betriebsordnung} zu
beachten, Werkzeug und Maschinen pfleglich zu behandeln, über Betriebs- und
Geschäftsgeheimnisse \emph{Stillschweigen} zu wahren und einen \emph{Ausbildungsnachweis}
(Berichtsheft) zu führen.

\subsection{Pflichten der Ausbildenden (\S\S~14--16 BBiG)}

Ausbildende haben dafür zu sorgen, dass die berufliche Handlungsfähigkeit vermittelt wird und
das Ausbildungsziel in der vorgesehenen Zeit erreicht werden kann. Sie müssen selbst
ausbilden oder eine:n Ausbilder:in beauftragen, die \emph{Ausbildungsmittel} (Werkzeuge,
Werkstoffe) kostenlos bereitstellen, zum Berufsschulbesuch anhalten und dafür sorgen, dass
Auszubildende nicht sittlich oder körperlich gefährdet werden (\emph{Fürsorgepflicht}). Nach
\S~15 BBiG sind Auszubildende für den \emph{Berufsschulunterricht} und für Prüfungen
freizustellen.

\begin{center}
\begin{tabular}{p{6.3cm} p{6.3cm}}
\toprule
\textbf{Pflichten der Auszubildenden (\S~13)} & \textbf{Pflichten der Ausbildenden (\S\S~14--16)}\\
\midrule
Aufgaben sorgfältig ausführen & Berufliche Handlungsfähigkeit vermitteln\\
Weisungen folgen, Ordnung beachten & Ausbildungsmittel kostenlos stellen\\
Berichtsheft führen & Zum Berufsschulbesuch anhalten, freistellen (\S~15)\\
Stillschweigen wahren & Fürsorge; angemessene Vergütung (\S~17)\\
Werkzeug pfleglich behandeln & Zeugnis ausstellen (\S~16)\\
\bottomrule
\end{tabular}
\end{center}

\subsection{Beendigung und Kündigung (\S~22 BBiG)}

Das Ausbildungsverhältnis endet regulär mit Ablauf der Ausbildungszeit oder --- bei Bestehen
der Prüfung --- mit deren Ergebnisbekanntgabe. Vorher kann es gekündigt werden:

\begin{itemize}
  \item \textbf{Während der Probezeit:} beide Seiten jederzeit, ohne Frist und ohne Angabe von
    Gründen (\S~22 Abs.~1).
  \item \textbf{Nach der Probezeit --- fristlos (außerordentlich):} beide Seiten nur aus einem
    \emph{wichtigen Grund}; die Kündigung muss innerhalb von zwei Wochen nach Bekanntwerden
    des Grundes erfolgen (\S~22 Abs.~2 u.~4).
  \item \textbf{Nach der Probezeit --- ordentlich (mit Frist von vier Wochen):} nur die/der
    \emph{Auszubildende}, wenn sie/er die Ausbildung aufgeben oder einen anderen Beruf
    anstreben will. Der Betrieb hat dieses ordentliche Kündigungsrecht nicht.
\end{itemize}

\begin{merksatz}
Der Gesetzgeber schützt das Ausbildungsverhältnis besonders: Nach der Probezeit kann der
\emph{Betrieb} nur noch aus wichtigem Grund kündigen. Ein ordentliches Kündigungsrecht (bloß
mit Frist) hat allein die/der Auszubildende.
\end{merksatz}

\subsection{Das Ausbildungszeugnis (\S~16 BBiG)}

Bei Beendigung ist der/dem Auszubildenden ein \emph{schriftliches} Zeugnis auszustellen (die
elektronische Form ist ausgeschlossen). Man unterscheidet:

\begin{itemize}
  \item \textbf{Einfaches Zeugnis:} Angaben über Art, Dauer und Ziel der Ausbildung sowie über
    die erworbenen Fertigkeiten, Kenntnisse und Fähigkeiten.
  \item \textbf{Qualifiziertes Zeugnis:} auf Verlangen der/des Auszubildenden zusätzlich
    Angaben über \emph{Verhalten und Leistung}.
\end{itemize}

\section{Jugendarbeitsschutzgesetz (JArbSchG)}

Für Auszubildende \emph{unter 18~Jahren} gilt neben dem BBiG das
\emph{Jugendarbeitsschutzgesetz}. Es definiert als \textbf{Kind}, wer noch nicht 15~Jahre alt
ist, und als \textbf{Jugendliche:n}, wer mindestens 15, aber noch nicht 18~Jahre alt ist
(\S~2 JArbSchG). Es schützt junge Menschen durch strengere Regeln:

\begin{center}
\begin{tabular}{p{4.8cm} p{7.8cm}}
\toprule
\textbf{Bereich} & \textbf{Kernregelung}\\
\midrule
Arbeitszeit & Höchstens 8~Std./Tag und 40~Std./Woche (\S~8)\\
Fünf-Tage-Woche & Beschäftigung in der Regel nur an fünf Tagen (\S~15)\\
Nachtruhe & Beschäftigung nur von 6 bis 20~Uhr (\S~14); danach mind.\ 12~Std.\ Freizeit (\S~13)\\
Ruhepausen & 30~Min.\ bei mehr als 4,5--6~Std., 60~Min.\ bei mehr als 6~Std.\ (\S~11)\\
Urlaub & Gestaffelt: mind.\ 30 / 27 / 25~Werktage je nach Alter (\S~19)\\
Freistellung & Für Berufsschule und Prüfungen (\S\S~9, 10)\\
Verbote & Keine gefährlichen Arbeiten, keine Akkordarbeit, keine Arbeit unter Tage\\
Untersuchungen & Erstuntersuchung vor Beginn (\S~32), Nachuntersuchung im 1.\ Jahr (\S~33)\\
\bottomrule
\end{tabular}
\end{center}

\begin{eselsbruecke}
\textbf{„6 bis 20, 5 Tage, 8 Stunden"} --- so merken Sie sich die drei wichtigsten Grenzen
für Jugendliche: Arbeit nur von 6 bis 20~Uhr, an fünf Tagen, höchstens acht Stunden am Tag.
\end{eselsbruecke}

Samstags- und Sonntagsarbeit sind grundsätzlich verboten (\S\S~16, 17 JArbSchG); für bestimmte
Branchen gelten Ausnahmen. Wichtig ist auch die Freistellung: An einem Berufsschultag mit mehr
als fünf Unterrichtsstunden darf die/der Jugendliche danach nicht mehr beschäftigt werden
(gilt für einen Tag pro Woche); beginnt der Unterricht vor 9~Uhr, ist eine Beschäftigung davor
unzulässig (\S~9 JArbSchG).

\section{Mitbestimmung: Betriebsrat und JAV}

Die Interessen der Arbeitnehmer werden durch das \emph{Betriebsverfassungsgesetz (BetrVG)}
geschützt. Zentrales Organ ist der \textbf{Betriebsrat}.

\subsection{Der Betriebsrat}

Ein Betriebsrat wird in Betrieben mit in der Regel mindestens \emph{fünf} ständigen
wahlberechtigten Arbeitnehmern gewählt, von denen drei wählbar sind (\S~1 BetrVG).
Wahlberechtigt (\emph{aktives Wahlrecht}) sind Arbeitnehmer ab 18~Jahren (\S~7); wählbar
(\emph{passives Wahlrecht}) ist, wer wahlberechtigt ist und dem Betrieb mindestens sechs
Monate angehört (\S~8). Die Zahl der Mitglieder richtet sich nach der Beschäftigtenzahl. Die
Amtszeit beträgt \emph{vier Jahre} (\S~21).

Der Betriebsrat wacht darüber, dass die Bestimmungen zugunsten der Arbeitnehmer eingehalten
werden, nimmt Anregungen entgegen und verhandelt mit dem Arbeitgeber. Seine Rechte sind
abgestuft: In \emph{sozialen} Angelegenheiten (z.\,B.\ Beginn/Ende der Arbeitszeit, Ordnung des
Betriebs) hat er echte \textbf{Mitbestimmungsrechte}; bei \emph{personellen} Einzelmaßnahmen
(Einstellung, Kündigung, Versetzung) hat er eingeschränkte Mitbestimmung; in
\emph{wirtschaftlichen} Angelegenheiten hat er Unterrichtungs- und Beratungsrechte.

\subsection{Die Jugend- und Auszubildendenvertretung (JAV)}

In Betrieben mit in der Regel mindestens fünf jugendlichen Arbeitnehmern (unter 18) oder
Auszubildenden (unter 25) wird eine \textbf{Jugend- und Auszubildendenvertretung} gewählt
(\S~60 BetrVG). Wählbar sind Beschäftigte, die das 25.\ Lebensjahr noch nicht vollendet haben;
Mitglieder des Betriebsrats können nicht in die JAV gewählt werden. Die Amtszeit beträgt
\emph{zwei Jahre} (\S~64).

Die JAV vertritt die Interessen der jungen Beschäftigten: Sie beantragt beim Betriebsrat
Maßnahmen (besonders zu Berufsbildung und Übernahme), wacht über die Einhaltung der für
Jugendliche und Auszubildende geltenden Vorschriften und fördert die Integration. Sie kann zu
Betriebsratssitzungen Vertreter entsenden.

\section{Zusammenfassung}

\begin{enumerate}
  \item Die \textbf{duale Ausbildung} findet an zwei Lernorten statt: Betrieb und Berufsschule; Rechtsrahmen ist das \textbf{BBiG}.
  \item Zu unterscheiden sind \textbf{Ausbildende:r} (Vertragspartner), \textbf{Ausbilder:in} (führt aus) und \textbf{Auszubildende:r}.
  \item Die Rechtsgrundlagen bauen aufeinander auf: BBiG $\rightarrow$ \textbf{Ausbildungsordnung} $\rightarrow$ \textbf{Ausbildungsrahmenplan} $\rightarrow$ \textbf{betrieblicher Ausbildungsplan}.
  \item Der \textbf{Berufsausbildungsvertrag} ist schriftlich abzuschließen, enthält neun Pflichtinhalte und wird bei der \textbf{IHK} eingetragen.
  \item \textbf{Probezeit} 1--4 Monate; \textbf{Vergütung} angemessen und jährlich steigend; \textbf{Urlaub} nach BUrlG bzw.\ JArbSchG.
  \item \textbf{Pflichten:} \S~13 (Auszubildende) und \S\S~14--16 (Ausbildende); \textbf{Kündigung} nach \S~22; \textbf{Zeugnis} einfach/qualifiziert (\S~16).
  \item Das \textbf{JArbSchG} schützt unter 18-Jährige (8~Std./Tag, 6--20~Uhr, 5-Tage-Woche, gestaffelter Urlaub, ärztliche Untersuchungen).
  \item \textbf{Betriebsrat} (ab 5 Arbeitnehmern, 4~Jahre) und \textbf{JAV} (unter 18/25, 2~Jahre) sichern die Mitbestimmung.
\end{enumerate}

\newpage

\section*{Glossar}
\addcontentsline{toc}{section}{Glossar}

\glossareintrag{Duale Ausbildung}{Berufsausbildung an zwei Lernorten --- Betrieb (Praxis) und Berufsschule (Theorie).}

\glossareintrag{Berufsbildungsgesetz (BBiG)}{Gesetz von 1969, das die Berufsausbildung in Deutschland regelt; zuletzt novelliert mit Wirkung ab 01.01.2020.}

\glossareintrag{Ausbildende:r}{Vertragspartner der/des Auszubildenden; das Unternehmen bzw.\ die Inhaberin/der Inhaber, das andere zur Berufsausbildung einstellt.}

\glossareintrag{Ausbilder:in}{Person, die die Berufsausbildung im Betrieb praktisch durchführt und von der/dem Ausbildenden damit beauftragt ist.}

\glossareintrag{Ausbildungsordnung}{Nach \S~5 BBiG für jeden Beruf erlassene Verordnung mit Berufsbild, Dauer, Ausbildungsrahmenplan und Prüfungsanforderungen.}

\glossareintrag{Ausbildungsrahmenplan}{Bundeseinheitliche grobe sachliche und zeitliche Gliederung der zu vermittelnden Inhalte; Teil der Ausbildungsordnung.}

\glossareintrag{Betrieblicher Ausbildungsplan}{Vom Betrieb erstellter, auf das Unternehmen abgestimmter Plan; Anlage des Berufsausbildungsvertrages.}

\glossareintrag{Berufsausbildungsvertrag}{Schriftlicher Vertrag zwischen Ausbildendem und Auszubildendem (\S\S~10, 11 BBiG); die elektronische Form ist ausgeschlossen.}

\glossareintrag{Zuständige Stelle}{Einrichtung, die das Ausbildungsverhältnis einträgt und überwacht; für kaufmännische Berufe die IHK.}

\glossareintrag{Probezeit}{Beginn des Ausbildungsverhältnisses; mindestens ein, höchstens vier Monate (\S~20 BBiG).}

\glossareintrag{Ausbildungsvergütung}{Angemessene Vergütung nach \S~17 BBiG; muss im Vertrag stehen und mindestens jährlich steigen.}

\glossareintrag{Ausbildungsnachweis}{Berichtsheft; von der/dem Auszubildenden zu führen (\S~13) und von der/dem Ausbildenden durchzusehen.}

\glossareintrag{Fürsorgepflicht}{Pflicht der/des Ausbildenden, dafür zu sorgen, dass Auszubildende charakterlich gefördert und nicht gefährdet werden.}

\glossareintrag{Außerordentliche Kündigung}{Fristlose Kündigung aus wichtigem Grund; innerhalb von zwei Wochen nach Bekanntwerden auszusprechen (\S~22 BBiG).}

\glossareintrag{Ordentliche Kündigung}{Kündigung mit vier Wochen Frist; nach der Probezeit nur durch die/den Auszubildende:n möglich (\S~22 BBiG).}

\glossareintrag{Ausbildungszeugnis}{Schriftliches Zeugnis bei Beendigung (\S~16 BBiG); einfach (Art/Dauer/Inhalte) oder qualifiziert (zusätzlich Verhalten und Leistung).}

\glossareintrag{Jugendarbeitsschutzgesetz (JArbSchG)}{Gesetz zum Schutz von unter 18-Jährigen; regelt Arbeitszeit, Pausen, Nachtruhe, Urlaub und Verbote.}

\glossareintrag{Betriebsverfassungsgesetz (BetrVG)}{Gesetz, das die betriebliche Interessenvertretung der Arbeitnehmer durch den Betriebsrat regelt.}

\glossareintrag{Betriebsrat}{Von den Arbeitnehmern gewähltes Organ der Mitbestimmung (ab 5 Arbeitnehmern, Amtszeit 4 Jahre).}

\glossareintrag{Jugend- und Auszubildendenvertretung (JAV)}{Vertretung der jugendlichen Arbeitnehmer (unter 18) und Auszubildenden (unter 25); Amtszeit 2 Jahre (\S\S~60 ff.\ BetrVG).}

\end{document}
